\chapter{Przestrzeń dualna i odwzorowania dualne}
\label{ch:lin7-przestrzen-dualna}
\index{przestrzeń dualna@Przestrzeń dualna}
\index{funkcjonał liniowy@Funkcjonał liniowy}
\index{anihilator@Anihilator}
\index{odwzorowanie dualne@Odwzorowanie dualne}
\index{bidualność@Bidualność}

\noindent\textbf{Cel rozdziału.}
Do tej pory badaliśmy wektory i odwzorowania
liniowe pomiędzy przestrzeniami wektorowymi.
Teraz każdemu wektorowi przyporządkujemy
możliwość oceniania go przez funkcjonały
liniowe, a każdemu odwzorowaniu liniowemu
--- odpowiadające mu odwzorowanie dualne.
Uzasadnimy konstrukcję baz dualnych,
anihilatorów i naturalne zanurzenie przestrzeni
w jej przestrzeń bidualną. Na końcu pokażemy,
jak z tego punktu widzenia rozumieć
ślad endomorfizmu.

W całym rozdziale \(K\) jest dowolnym ciałem,
\(0_K\ne1_K\), a \(V,W,Z\) oznaczają przestrzenie
wektorowe nad \(K\). Nie zakładamy, że
charakterystyka \(K\) jest różna od \(2\).
Z wyjątkiem wyraźnie oznaczonych uwag
o wymiarze nieskończonym wszystkie twierdzenia
dotyczą przestrzeni skończenie wymiarowych;
dopuszczamy także wymiar \(0\).
W razie potrzeby zapis \(\operatorname{Hom}_K(V,W)\)
ma znaczenie ustalone w
Rozdziale~\ref{ch:lin3-odwzorowania-liniowe}.

\section{Funkcjonały liniowe i przestrzeń dualna}
\label{sec:lin7-funkcjonaly}

\begin{definicja}[Funkcjonał i przestrzeń dualna]
\label{def:lin7-dual}
\emph{Funkcjonałem liniowym} na \(V\)
nazywamy odwzorowanie liniowe
\(\varphi\colon V\to K\), gdzie \(K\)
traktujemy jako jednowymiarową
przestrzeń wektorową nad sobą.
\emph{Przestrzeń dualna} do \(V\) to
\[
 V^*:=\operatorname{Hom}_K(V,K).
\]
Dodawanie i mnożenie funkcjonałów
przez skalary określamy punktowo:
\((\varphi+\psi)(v)=\varphi(v)+\psi(v)\)
oraz \((a\varphi)(v)=a\varphi(v)\).
Parę \((\varphi,v)\in V^*\times V\)
łączymy zapisem
\(\langle\varphi,v\rangle:=\varphi(v)\).
\end{definicja}

\begin{stwierdzenie}[Liniowość ewaluacji]
\label{prop:lin7-ewaluacja}
Dla dowolnych \(\varphi,\psi\in V^*\),
\(v,w\in V\) i \(a,b\in K\) zachodzi
\[
 \begin{aligned}
 \langle a\varphi+b\psi,v\rangle
   &=a\langle\varphi,v\rangle
     +b\langle\psi,v\rangle,\\
 \langle\varphi,av+bw\rangle
   &=a\langle\varphi,v\rangle
     +b\langle\varphi,w\rangle.
 \end{aligned}
\]
Zatem ewaluacja jest funkcją
dwuliniową \(V^*\times V\to K\).
Jeżeli \(\dim V<\infty\) i \(v\ne0\),
to istnieje \(\varphi\in V^*\)
o \(\varphi(v)\ne0_K\).
\end{stwierdzenie}
\begin{proof}
Pierwsza równość wynika z punktowego
określenia działań w \(V^*\),
a druga z liniowości \(\varphi\).
Dla \(v\ne0\) jednoelementowy zbiór
\(\{v\}\) jest liniowo niezależny.
Uzupełnijmy go do bazy
\((v,v_2,\ldots,v_n)\) przestrzeni \(V\),
co jest możliwe na mocy
Twierdzenia~\ref{thm:lin-uzupelnianie-skonczone}.
Przyporządkowanie
\(v\mapsto1_K\), \(v_i\mapsto0_K\)
dla \(i\geq2\) przedłuża się
jednoznacznie do funkcjonału liniowego
na mocy
Twierdzenia~\ref{thm:lin3-przedluzenie-z-bazy}.
Dla niego \(\varphi(v)=1_K\ne0_K\).
\end{proof}

\begin{przyklad}[Funkcjonały współrzędnych]
\label{prz:lin7-wspolrzedne}
Niech \(V=K^3\). Wzory
\[
 \varphi(x,y,z)=2x-y+3z,\qquad
 \psi(x,y,z)=x+z
\]
określają funkcjonały liniowe.
Każdy funkcjonał
\(\lambda\colon K^3\to K\)
ma postać
\(\lambda(x,y,z)=ax+by+cz\)
dla dokładnie jednej trójki
\((a,b,c)\in K^3\), ponieważ
wartości na wektorach standardowych
wyznaczają go jednoznacznie.
Współczynniki \(2,-1,3\)
interpretujemy w dowolnym ciele \(K\);
w szczególności mogą się one
upraszczać w dodatniej charakterystyce.
\end{przyklad}

\begin{uwaga}[Algebraiczna przestrzeń dualna nie wymaga topologii]
\label{rem:lin7-dual-algebraiczny}
W definicji \(V^*\) nie występują
norma, iloczyn skalarny ani topologia.
Wszystkie funkcjonały są liniowe,
ale nie narzucamy na nie dodatkowego
warunku ciągłości. Jeśli w późniejszych
rozdziałach pojawi się norma, przestrzeń
ciągłych funkcjonałów liniowych będzie
wymagała osobnej definicji.
\end{uwaga}

\section{Bazy dualne i ich współrzędne}
\label{sec:lin7-baza-dualna}

\begin{definicja}[Baza dualna]
\label{def:lin7-baza-dualna}
Niech \(\dim V=n<\infty\) i
\(B=(v_1,\ldots,v_n)\) będzie
uporządkowaną bazą \(V\).
Dla każdego \(i\in\{1,\ldots,n\}\)
\emph{funkcjonałem dualnym}
do wektora \(v_i\) nazywamy
jedyny funkcjonał \(v^i\in V^*\)
spełniający
\[
 v^i(v_j)=\delta_{ij}
 \qquad (1\leq j\leq n).
\]
Układ \(B^*=(v^1,\ldots,v^n)\)
nazywamy \emph{bazą dualną} do \(B\).
W przypadku \(n=0\) obie bazy
są pustymi układami.
\end{definicja}

\begin{twierdzenie}[Istnienie i własność bazy dualnej]
\label{thm:lin7-baza-dualna}
Niech \(V\) będzie przestrzenią nad \(K\)
o wymiarze \(n<\infty\), a
\(B=(v_1,\ldots,v_n)\) jej bazą.
Wówczas funkcjonały
\(v^1,\ldots,v^n\) istnieją,
są jednoznaczne i tworzą bazę \(V^*\).
Dla każdych \(v\in V\) i
\(\varphi\in V^*\) zachodzą
tożsamości
\[
 v=\sum_{i=1}^{n}v^i(v)v_i,
 \qquad
 \varphi=\sum_{i=1}^{n}\varphi(v_i)v^i.
\]
W szczególności \(\dim V^*=\dim V=n\).
\end{twierdzenie}
\begin{proof}
Z
Twierdzenia~\ref{thm:lin3-przedluzenie-z-bazy}
wynika, że każde przyporządkowanie
skalarów wektorom bazy \(B\)
określa dokładnie jeden funkcjonał
liniowy. Dla każdego \(i\) wybierzmy
wartości \(\delta_{ij}\), co dowodzi
istnienia i jednoznaczności \(v^i\).

Niech \(v=\sum_j a_jv_j\).
Stosując \(v^i\), otrzymujemy
\(v^i(v)=\sum_j a_j\delta_{ij}=a_i\);
stąd pierwsza tożsamość.
Niech \(\varphi\in V^*\) i określmy
\(\psi=\sum_i\varphi(v_i)v^i\).
Dla każdego \(j\) mamy
\(\psi(v_j)=\sum_i\varphi(v_i)\delta_{ij}
=\varphi(v_j)\).
Dwa funkcjonały zgodne na bazie są
równe, więc \(\psi=\varphi\).
Oznacza to, że \(B^*\) rozpina \(V^*\).

Jeśli \(\sum_i a_iv^i=0\),
to po obliczeniu wartości w \(v_j\)
otrzymujemy \(a_j=0_K\) dla każdego
\(j\). Układ \(B^*\) jest więc
liniowo niezależny.
Dla \(n=0\) mamy \(V=\{0\}\);
jedyny funkcjonał z \(V\) do \(K\)
jest zerowy, a puste sumy
dają wymagane tożsamości.
\end{proof}

\begin{wniosek}[Współrzędne funkcjonału]
\label{cor:lin7-wspolrzedne-dualne}
Przy założeniach
Twierdzenia~\ref{thm:lin7-baza-dualna}
wektor kolumnowy współrzędnych
\(\varphi\in V^*\) w bazie dualnej
jest równy
\[
 [\varphi]_{B^*}
 =\bigl(\varphi(v_1),\ldots,\varphi(v_n)
   \bigr)^{\mathsf T}\in K^n.
\]
Jeżeli \([v]_B=(x_1,\ldots,x_n)^{\mathsf T}\),
to
\[
 \varphi(v)=[\varphi]_{B^*}^{\mathsf T}[v]_B
 =\sum_{i=1}^n \varphi(v_i)x_i.
\]
\end{wniosek}
\begin{proof}
Pierwsze twierdzenie daje rozwinięcie
\(\varphi=\sum_i\varphi(v_i)v^i\),
czyli wskazane współrzędne.
Ponieważ \(v=\sum_i x_iv_i\),
liniowość \(\varphi\) daje drugą równość.
Dla \(n=0\) suma jest pusta i wynosi
\(0_K\).
\end{proof}

\begin{przyklad}[Baza dualna w \(\mathbb R^2\)]
\label{prz:lin7-baza-dualna-r2}
Niech \(V=\mathbb R^2\) i
\(B=(v_1,v_2)\), gdzie
\(v_1=(1,1)^{\mathsf T}\) i
\(v_2=(1,-1)^{\mathsf T}\).
Dla \(x=(s,t)^{\mathsf T}\) rozwiązanie
\(x=av_1+bv_2\) daje
\(a=(s+t)/2\), \(b=(s-t)/2\).
Zatem
\[
 v^1(s,t)=\frac{s+t}{2},\qquad
 v^2(s,t)=\frac{s-t}{2}.
\]
Istotnie \(v^1(v_1)=v^2(v_2)=1\)
i \(v^1(v_2)=v^2(v_1)=0\).
Wzory wykorzystują dzielenie przez
\(2\), dlatego ten konkretny przykład
dotyczy \(\mathbb R\), a nie ciała
charakterystyki \(2\).
\end{przyklad}

\begin{twierdzenie}[Zmiana bazy dualnej]
\label{thm:lin7-zmiana-bazy-dualnej}
Niech \(V\) ma wymiar \(n<\infty\),
a \(B=(v_1,\ldots,v_n)\) i
\(C=(w_1,\ldots,w_n)\) będą bazami
uporządkowanymi. Niech \(P\in
\operatorname{Mat}_{n\times n}(K)\)
będzie macierzą przejścia określoną
równością
\[
 [v]_B=P[v]_C\qquad(v\in V).
\]
Wtedy dla każdego \(\varphi\in V^*\)
zachodzi
\[
 [\varphi]_{C^*}=P^{\mathsf T}
 [\varphi]_{B^*},
\]
a funkcjonały bazowe spełniają
\[
 w^j=\sum_{i=1}^n (P^{-1})_{ji}v^i,
 \qquad 1\leq j\leq n.
\]
\end{twierdzenie}
\begin{proof}
Kolumna \(j\) macierzy \(P\)
zawiera współrzędne \(w_j\) w bazie
\(B\), więc
\(w_j=\sum_i p_{ij}v_i\).
Dla każdego \(\varphi\in V^*\)
otrzymujemy
\(\varphi(w_j)=\sum_i p_{ij}\varphi(v_i)\).
Z
Wniosku~\ref{cor:lin7-wspolrzedne-dualne}
wynika, że jest to dokładnie
\(j\)-ta współrzędna wektora
\(P^{\mathsf T}[\varphi]_{B^*}\).
Dowodzi to pierwszego wzoru.

Macierz współrzędnych wartości
\((v^i(w_j))_{i,j}\) jest równa \(P\).
Pisząc \(w^j=\sum_i c_{ji}v^i\)
i podstawiając \(w_k\), otrzymujemy
warunek
\(\sum_i c_{ji}p_{ik}=\delta_{jk}\).
Zatem macierz \(C=(c_{ji})\)
spełnia \(CP=I_n\), czyli
\(C=P^{-1}\). Daje to drugi wzór.
Dla \(n=0\) macierz przejścia to
\(I_0\), a wszystkie równości są puste.
\end{proof}

\begin{uwaga}[Brak kanonicznej identyfikacji \(V\) i \(V^*\)]
\label{rem:lin7-niekanoniczna}
Równość wymiarów gwarantuje
istnienie izomorfizmu \(V\cong V^*\),
lecz powyższy wybór
\(v_i\mapsto v^i\) zależy od bazy.
Przy jej zmianie wektory
i funkcjonały mają odmienne reguły
przekształcania współrzędnych.
Nie należy więc utożsamiać
\(V\) z \(V^*\) bez dodatkowej struktury.
\end{uwaga}

\section{Anihilatory podprzestrzeni}
\label{sec:lin7-anihilatory}

\begin{definicja}[Anihilator]
\label{def:lin7-anihilator}
Niech \(U\) będzie podprzestrzenią
przestrzeni wektorowej \(V\).
Jej \emph{anihilatorem} nazywamy
\[
 U^\circ:=
 \{\varphi\in V^*:\ \varphi(u)=0_K
                 \text{ dla każdego }u\in U\}.
\]
Jest to podprzestrzeń \(V^*\).
Dla podprzestrzeni \(S\subset V^*\)
oznaczamy również przez
\[
 S_\circ:=
 \{v\in V:\ \varphi(v)=0_K
                   \text{ dla każdego }\varphi\in S\}
\]
jej anihilator w \(V\).
\end{definicja}

\begin{stwierdzenie}[Monotoniczność i przypadki skrajne]
\label{prop:lin7-anihilator-wlasnosci}
Niech \(V\) będzie skończenie wymiarowa,
a \(U_1,U_2\subset V\) podprzestrzeniami.
Wtedy
\[
 U_1\subset U_2
 \quad\Longrightarrow\quad
 U_2^\circ\subset U_1^\circ,
 \qquad
 \{0\}^\circ=V^*,\quad
 V^\circ=\{0\}.
\]
Dla dowolnych podprzestrzeni \(U_1,U_2\)
zachodzi ponadto
\[
 (U_1+U_2)^\circ=U_1^\circ\cap U_2^\circ.
\]
\end{stwierdzenie}
\begin{proof}
Funkcjonał zerujący się na większej
podprzestrzeni zeruje się na mniejszej,
co daje pierwszą inkluzję.
Każdy funkcjonał zeruje się w \(0\),
zaś funkcjonał zerujący się na całym
\(V\) jest funkcjonałem zerowym.
Wreszcie
\(\varphi(u_1+u_2)=0_K\)
dla wszystkich \(u_1\in U_1\),
\(u_2\in U_2\) wtedy i tylko wtedy,
gdy \(\varphi\) zeruje się na obu
podprzestrzeniach: dla implikacji
w prawo podstawiamy kolejno
\(u_2=0\) i \(u_1=0\);
dla implikacji odwrotnej korzystamy
z liniowości \(\varphi\).
\end{proof}

\begin{twierdzenie}[Wymiar anihilatora]
\label{thm:lin7-wymiar-anihilatora}
Niech \(V\) ma wymiar \(n<\infty\),
a \(U\subset V\) będzie
podprzestrzenią o wymiarze \(r\).
Wtedy
\[
 \dim U^\circ=n-r.
\]
Co więcej, jeśli
\(B=(u_1,\ldots,u_r,v_{r+1},\ldots,v_n)\)
jest bazą \(V\), w której pierwsze
\(r\) wektorów tworzy bazę \(U\),
to
\[
 U^\circ=
 \operatorname{span}\{v^{r+1},\ldots,v^n\},
\]
gdzie \(v^{r+1},\ldots,v^n\)
są odpowiednimi elementami
bazy dualnej \(B^*\).
\end{twierdzenie}
\begin{proof}
Dzięki twierdzeniu o uzupełnianiu
bazy możemy wybrać bazę \(B\)
o wskazanej postaci. Każdy funkcjonał
\(\varphi\in V^*\) zapisujemy
jednoznacznie jako
\[
 \varphi=\sum_{i=1}^{r}
          \varphi(u_i)u^i
        +\sum_{i=r+1}^{n}
          \varphi(v_i)v^i.
\]
Funkcjonał zeruje się na \(U\)
wtedy i tylko wtedy, gdy
\(\varphi(u_i)=0_K\) dla wszystkich
\(1\leq i\leq r\), ponieważ wektory
\(u_1,\ldots,u_r\) rozpinają \(U\).
Zatem pozostaje druga suma.
Odpowiednie funkcjonały tworzą
liniowo niezależną rodzinę
długości \(n-r\) i generują
\(U^\circ\). To daje obie tezy.
Obejmuje to \(U=\{0\}\),
\(U=V\), a także \(n=0\).
\end{proof}

\begin{wniosek}[Anihilator anihilatora]
\label{cor:lin7-podwojny-anihilator}
Jeżeli \(V\) ma skończony wymiar,
a \(U\subset V\) i \(S\subset V^*\)
są podprzestrzeniami, to
\[
 (U^\circ)_\circ=U,\qquad
 (S_\circ)^\circ=S.
\]
\end{wniosek}
\begin{proof}
Przyjmijmy bazę \(B\) jak
w Twierdzeniu~\ref{thm:lin7-wymiar-anihilatora}.
Wtedy \(U^\circ\) jest rozpięty
przez \(v^{r+1},\ldots,v^n\).
Wektor \(x=\sum_{i=1}^n x_iv_i\)
należy do \((U^\circ)_\circ\)
wtedy i tylko wtedy, gdy
\(x_{r+1}=\cdots=x_n=0_K\).
Oznacza to dokładnie, że \(x\in U\).

Dla drugiej równości wybierzmy bazę
\((\varphi_1,\ldots,\varphi_s)\)
przestrzeni \(S\), gdzie \(s=\dim S\).
Określmy odwzorowanie liniowe
\[
 L\colon V\to K^s,\qquad
 L(v)=(\varphi_1(v),\ldots,\varphi_s(v)).
\]
Jest ono surjektywne. Istotnie,
gdyby \(\operatorname{im}L\ne K^s\),
istniałby niezerowy funkcjonał
\(\eta\in(K^s)^*\) zerujący się
na \(\operatorname{im}L\):
wybieramy bazę obrazu, uzupełniamy
ją do bazy \(K^s\) i nadajemy
wartość \(1_K\) jednemu z
uzupełniających wektorów,
a pozostałym wartość \(0_K\).
W bazie standardowej funkcjonał
\(\eta\) ma pewne współczynniki
\(a_1,\ldots,a_s\), nie wszystkie zerowe.
Jednak \(\eta\circ L=0\)
oznaczałoby
\(\sum_{i=1}^s a_i\varphi_i=0\),
wbrew niezależności liniowej bazy \(S\).
Zatem \(\operatorname{rk}L=s\).

Z definicji \(\ker L=S_\circ\).
Twierdzenie o wymiarze jądra
i obrazu daje
\(\dim S_\circ=n-s\), a
Twierdzenie~\ref{thm:lin7-wymiar-anihilatora}
zapewnia
\(\dim(S_\circ)^\circ=s\).
Każdy funkcjonał z \(S\) zeruje się
na \(S_\circ\), czyli
\(S\subset(S_\circ)^\circ\).
Inkluzja między podprzestrzeniami
tego samego wymiaru jest równością.
W przypadku \(s=0\) odwzorowanie
\(L\colon V\to K^0\) jest surjektywne
i ten sam argument daje wynik.
\end{proof}

\begin{przyklad}[Anihilator płaszczyzny]
\label{prz:lin7-anihilator-r3}
Niech \(V=\mathbb R^3\), a
\[
 U=\{(x,y,z)\in\mathbb R^3:
                 x+2y-z=0\}.
\]
Przestrzeń \(U\) ma wymiar \(2\).
Funkcjonał \(\ell(x,y,z)=x+2y-z\)
zeruje się na \(U\), więc
\(\ell\in U^\circ\).
Ponieważ \(\ell\ne0\), a
\(\dim U^\circ=3-2=1\),
otrzymujemy
\(U^\circ=\operatorname{span}\{\ell\}\).
Przykład pokazuje, że jedno równanie
liniowe opisujące podprzestrzeń
jest w istocie elementem przestrzeni dualnej.
\end{przyklad}

\begin{twierdzenie}[Anihilator przecięcia]
\label{thm:lin7-anihilator-przeciecia}
Jeśli \(V\) ma skończony wymiar
i \(U_1,U_2\subset V\) są
podprzestrzeniami, to
\[
 (U_1\cap U_2)^\circ
 =U_1^\circ+U_2^\circ.
\]
\end{twierdzenie}
\begin{proof}
Jeśli
\(\varphi=\varphi_1+\varphi_2\),
gdzie \(\varphi_i\in U_i^\circ\),
to \(\varphi\) zeruje się na
\(U_1\cap U_2\), więc
\(U_1^\circ+U_2^\circ
 \subset (U_1\cap U_2)^\circ\).
Dla równości wymiarów
korzystamy z
Twierdzenia~\ref{thm:lin2-wzor-wymiarowy},
z równości
\((U_1+U_2)^\circ
=U_1^\circ\cap U_2^\circ\)
i ze wzoru na wymiar anihilatora.
Mamy
\[
\begin{aligned}
 \dim(U_1^\circ+U_2^\circ)
 &= (n-\dim U_1)+(n-\dim U_2)\\
 &\quad -(n-\dim(U_1+U_2))\\
 &=n-\dim(U_1\cap U_2)\\
 &=\dim(U_1\cap U_2)^\circ.
\end{aligned}
\]
Inkluzja podprzestrzeni o tych
samych skończonych wymiarach
jest równością.
\end{proof}

\section{Przestrzeń dualna do ilorazu i rozszerzanie funkcjonałów}
\label{sec:lin7-dual-ilorazu}

\begin{twierdzenie}[Rozszerzenie funkcjonału z podprzestrzeni]
\label{thm:lin7-rozszerzanie}
Niech \(V\) ma skończony wymiar,
\(U\subset V\) będzie podprzestrzenią,
a \(\lambda\in U^*\).
Wówczas istnieje
\(\widetilde\lambda\in V^*\)
spełniający
\(\widetilde\lambda(u)=\lambda(u)\)
dla każdego \(u\in U\).
Nie twierdzimy, że to rozszerzenie
jest jednoznaczne.
\end{twierdzenie}
\begin{proof}
Wybierzmy bazę
\((u_1,\ldots,u_r)\) przestrzeni \(U\).
Uzupełnijmy ją do bazy \(V\)
o pozostałych wektorach
\(v_{r+1},\ldots,v_n\).
Wartości \(\lambda(u_i)\)
dla \(1\leq i\leq r\)
zachowujemy, a na wektorach
\(v_{r+1},\ldots,v_n\)
ustalamy wartości \(0_K\).
Twierdzenie o odwzorowaniu zadanym
na bazie daje jednoznaczny funkcjonał
na \(V\) o tych wartościach.
Jego ograniczenie do \(U\)
jest równe \(\lambda\),
gdyż oba funkcjonały są liniowe
i zgodne na bazie \(U\).
Dla \(U=\{0\}\) jedyny funkcjonał
na \(U\) jest zerowy i wystarcza
wybrać \(\widetilde\lambda=0\).
\end{proof}

\begin{twierdzenie}[Przestrzeń dualna do ilorazu]
\label{thm:lin7-dual-ilorazu}
Niech \(V\) będzie skończenie wymiarowa,
\(U\subset V\) podprzestrzenią,
a \(q\colon V\to V/U\),
\(q(v)=v+U\), odwzorowaniem ilorazowym.
Odwzorowanie
\[
 q^*\colon (V/U)^*\longrightarrow V^*,
 \qquad q^*(\psi)=\psi\circ q,
\]
jest liniowe i różnowartościowe,
a jego obrazem jest dokładnie
\(U^\circ\).
W konsekwencji
\[
 (V/U)^*\cong U^\circ.
\]
\end{twierdzenie}
\begin{proof}
Złożenie dwóch odwzorowań liniowych
jest liniowe, więc \(q^*\)
jest dobrze określone.
Jego liniowość wynika
z punktowych działań na funkcjonałach.
Jeśli \(\psi\circ q=0\),
to z surjektywności \(q\)
wynika \(\psi=0\), więc \(q^*\)
jest różnowartościowe.
Jeśli \(\psi\in(V/U)^*\),
to dla \(u\in U\) mamy \(q(u)=0\)
oraz \((\psi\circ q)(u)=0_K\).
Stąd \(\operatorname{im}q^*\subset U^\circ\).

Odwrotnie, jeśli
\(\varphi\in U^\circ\),
to wzór
\(\psi(v+U):=\varphi(v)\)
jest dobrze określony:
gdy \(v+U=w+U\), wówczas
\(v-w\in U\), a więc
\(\varphi(v-w)=0_K\).
Wzór określa funkcjonał liniowy
na \(V/U\), dla którego
\(\psi\circ q=\varphi\).
Zatem \(\operatorname{im}q^*=U^\circ\).
\end{proof}

\begin{wniosek}[Przestrzeń dualna do sumy prostej]
\label{cor:lin7-dual-sumy-prostej}
Niech \(V=U_1\oplus U_2\)
będzie wewnętrzną sumą prostą
skończenie wymiarowych podprzestrzeni.
Wtedy odwzorowanie
\[
 V^*\longrightarrow U_1^*\oplus U_2^*,
 \qquad
 \varphi\longmapsto
   (\varphi|_{U_1},\varphi|_{U_2})
\]
jest izomorfizmem przestrzeni
wektorowych nad \(K\).
\end{wniosek}
\begin{proof}
Odwzorowanie jest liniowe.
Jeśli oba ograniczenia funkcjonału
są zerowe, to dla dowolnego
\(v=u_1+u_2\) mamy
\(\varphi(v)=0_K+0_K=0_K\);
odwzorowanie jest więc iniektywne.
Dla dowolnych
\(\lambda_i\in U_i^*\)
wzór
\(\varphi(u_1+u_2)=
\lambda_1(u_1)+\lambda_2(u_2)\)
jest dobrze określony, ponieważ
rozkład w sumie prostej jest
jednoznaczny. Jest liniowy i
ma wymagane ograniczenia, zatem
odwzorowanie jest surjektywne.
\end{proof}

\section{Odwzorowanie dualne i macierz transponowana}
\label{sec:lin7-odwzorowanie-dualne}

\begin{definicja}[Odwzorowanie dualne]
\label{def:lin7-odwzorowanie-dualne}
Dla dowolnego odwzorowania liniowego
\(T\colon V\to W\)
\emph{odwzorowaniem dualnym} do \(T\)
nazywamy
\[
 T^*\colon W^*\to V^*,
 \qquad
 (T^*\psi)(v):=\psi(Tv),
 \qquad \psi\in W^*,\ v\in V.
\]
A więc \(T^*(\psi)=\psi\circ T\).
Kierunek działania jest przeciwny
do kierunku \(T\).
\end{definicja}

\begin{twierdzenie}[Liniowość i odwracanie składania]
\label{thm:lin7-funktorialnosc}
Niech \(T\colon V\to W\)
i \(S\colon W\to Z\) będą
odwzorowaniami liniowymi.
Wówczas \(T^*\colon W^*\to V^*\)
jest liniowe oraz
\[
 (S\circ T)^*=T^*\circ S^*,
 \qquad
 (\operatorname{id}_V)^*
 =\operatorname{id}_{V^*}.
\]
Jeżeli \(T\) jest izomorfizmem,
to \(T^*\) jest izomorfizmem oraz
\[
 (T^*)^{-1}=(T^{-1})^*.
\]
\end{twierdzenie}
\begin{proof}
Dla \(\psi_1,\psi_2\in W^*\),
\(a,b\in K\) i \(v\in V\) mamy
\[
 \bigl(T^*(a\psi_1+b\psi_2)\bigr)(v)
 =(a\psi_1+b\psi_2)(Tv)
 =a(T^*\psi_1)(v)+b(T^*\psi_2)(v).
\]
Ponieważ funkcjonały o identycznych
wartościach na wszystkich wektorach
są równe, dowodzi to liniowości \(T^*\).

Dla \(\eta\in Z^*\) i \(v\in V\)
obliczamy
\[
 \bigl((S\circ T)^*\eta\bigr)(v)
 =\eta(S(Tv))
 =\bigl(T^*(S^*\eta)\bigr)(v),
\]
co daje równość złożeń.
Wzór dla identyczności wynika
bezpośrednio z definicji.

Jeśli \(T^{-1}\) istnieje,
to odwracanie składania daje
\[
 T^*\circ(T^{-1})^*
 =(T^{-1}\circ T)^*
 =\operatorname{id}_{V^*}
\]
oraz
\((T^{-1})^*\circ T^*
 =\operatorname{id}_{W^*}\).
Zatem \((T^{-1})^*\)
jest dwustronną odwrotnością \(T^*\).
\end{proof}

\begin{twierdzenie}[Macierz odwzorowania dualnego]
\label{thm:lin7-macierz-dualnego}
Niech \(\dim V=n<\infty\) i
\(\dim W=m<\infty\).
Wybierzmy uporządkowane bazy
\(B=(v_1,\ldots,v_n)\) w \(V\)
oraz \(C=(w_1,\ldots,w_m)\) w \(W\).
Jeżeli \(A=(a_{ij})\in
\operatorname{Mat}_{m\times n}(K)\)
jest macierzą odwzorowania liniowego
\(T\colon V\to W\) w bazach \(B,C\),
to macierz odwzorowania
\(T^*\colon W^*\to V^*\)
w bazach dualnych \(C^*,B^*\)
jest równa \(A^{\mathsf T}\).
W szczególności
\[
 [T^*\psi]_{B^*}
  =A^{\mathsf T}[\psi]_{C^*}
 \qquad(\psi\in W^*).
\]
\end{twierdzenie}
\begin{proof}
Z definicji macierzy odwzorowania
\(Tv_j=\sum_{i=1}^m a_{ij}w_i\).
Dla funkcjonału \(w^i\)
z bazy dualnej \(C^*\)
i \(v_j\in B\) mamy
\[
 (T^*w^i)(v_j)
 =w^i(Tv_j)
 =w^i\Bigl(\sum_{k=1}^m a_{kj}w_k\Bigr)
 =a_{ij}.
\]
Zatem \(j\)-ta współrzędna
funkcjonału \(T^*w^i\)
w bazie \(B^*\) jest równa \(a_{ij}\).
Jest to element w wierszu \(j\)
i kolumnie \(i\) macierzy \(T^*\),
czyli element \((A^{\mathsf T})_{ji}\).
Gdy \(n=0\) lub \(m=0\),
obliczenie dotyczy pustych rodzin
i nadal daje właściwą macierz
rozmiaru \(n\times m\).
\end{proof}

\begin{figure}[!htbp]
\centering
\begin{tikzpicture}[>=Latex,node distance=33mm and 37mm]
 \node (wstar) {\(W^*\)};
 \node[right=of wstar] (vstar) {\(V^*\)};
 \node[below=of wstar] (km) {\(K^m\)};
 \node[right=of km] (kn) {\(K^n\)};
 \draw[->] (wstar)--node[above]{\(T^*\)}(vstar);
 \draw[->] (wstar)--node[left]{\([\cdot]_{C^*}\)}(km);
 \draw[->] (vstar)--node[right]{\([\cdot]_{B^*}\)}(kn);
 \draw[->] (km)--node[below]{\(A^{\mathsf T}\)}(kn);
\end{tikzpicture}
\caption{Odwzorowanie dualne zmienia kierunek
względem \(T\colon V\to W\).
W bazach dualnych jego macierzą jest
transpozycja macierzy \(T\).}
\label{fig:lin7-dualne-diagram}
\end{figure}

\begin{przyklad}[Działanie odwzorowania dualnego]
\label{prz:lin7-transpozycja}
Niech \(V=\mathbb R^2\), \(W=\mathbb R^3\)
i niech w bazach standardowych
odwzorowanie \(T\) ma macierz
\[
 A=\begin{pmatrix}
 1&2\\
 -1&0\\
 3&1
 \end{pmatrix}.
\]
Dla \(\psi(x,y,z)=2x-y+z\)
kolumna współrzędnych \(\psi\)
w bazie dualnej wynosi
\((2,-1,1)^{\mathsf T}\).
Zatem
\[
 [T^*\psi]_{B^*}
 =\begin{pmatrix}1&-1&3\\2&0&1\end{pmatrix}
  \begin{pmatrix}2\\-1\\1\end{pmatrix}
 =\begin{pmatrix}6\\5\end{pmatrix}.
\]
Bez użycia macierzy:
\(T(s,t)=(s+2t,-s,3s+t)\),
więc
\(\psi(T(s,t))=2(s+2t)+s+(3s+t)=6s+5t\).
Obie metody dają ten sam funkcjonał.
\end{przyklad}

\section{Jądro, obraz i rząd odwzorowania dualnego}
\label{sec:lin7-dual-rzad}

\begin{twierdzenie}[Jądro i obraz odwzorowania dualnego]
\label{thm:lin7-jadro-obraz-dualnego}
Niech \(V,W\) będą
skończenie wymiarowymi przestrzeniami
nad \(K\), a \(T\colon V\to W\)
odwzorowaniem liniowym. Wtedy
\[
 \ker T^*=(\operatorname{im}T)^\circ
 \quad\text{(podprzestrzeń \(W^*\))},
 \qquad
 \operatorname{im}T^*=(\ker T)^\circ
 \quad\text{(podprzestrzeń \(V^*\))}.
\]
W szczególności
\[
 \operatorname{rk}T^*=\operatorname{rk}T.
\]
\end{twierdzenie}
\begin{proof}
Dla \(\psi\in W^*\) warunek
\(\psi\in\ker T^*\)
jest równoważny
\((T^*\psi)(v)=\psi(Tv)=0_K\)
dla każdego \(v\in V\).
Oznacza to dokładnie
\(\psi|_{\operatorname{im}T}=0\),
czyli pierwszą równość.

Jeśli \(\psi\in W^*\) i \(v\in\ker T\),
to \((T^*\psi)(v)=\psi(0)=0_K\).
Daje to inkluzję
\(\operatorname{im}T^*
 \subset(\ker T)^\circ\).

Dla drugiej inkluzji niech
\(\varphi\in(\ker T)^\circ\).
Na \(\operatorname{im}T\subset W\)
określmy funkcjonał
\(\lambda(Tv)=\varphi(v)\).
Jeżeli \(Tv=Tw\), to \(v-w\in\ker T\),
a więc \(\varphi(v-w)=0_K\);
wzór jest zatem dobrze określony.
Jest liniowy, ponieważ
\(\lambda(T(av+bw))
=\varphi(av+bw)
=a\varphi(v)+b\varphi(w)\).
Na mocy
Twierdzenia~\ref{thm:lin7-rozszerzanie}
istnieje funkcjonał \(\psi\in W^*\)
rozszerzający \(\lambda\).
Dla każdego \(v\in V\) otrzymujemy
\((T^*\psi)(v)=\psi(Tv)=\varphi(v)\),
zatem \(\varphi=T^*\psi\).
Wykazaliśmy równość obrazów.

Z twierdzenia o wymiarze jądra
i obrazu oraz wzoru na wymiar
anihilatora wynika
\[
 \begin{aligned}
 \operatorname{rk}T^*
 &=\dim(\ker T)^\circ\\
 &=\dim V-\dim\ker T\\
 &=\operatorname{rk}T.
 \end{aligned}
\]
Rachunek obejmuje również
przypadki \(V=\{0\}\) i \(W=\{0\}\).
\end{proof}

\begin{wniosek}[Rząd macierzy i jej transpozycji]
\label{cor:lin7-rzad-transpozycji}
Dla dowolnych \(m,n\in\mathbb N_0\)
i macierzy \(A\in
\operatorname{Mat}_{m\times n}(K)\)
zachodzi
\[
 \operatorname{rk}A^{\mathsf T}
 =\operatorname{rk}A.
\]
Jest to równość rzędów kolumnowych,
uzyskana bez używania iloczynu
skalarnego na \(K^m\) lub \(K^n\).
\end{wniosek}
\begin{proof}
Niech \(T\colon K^n\to K^m\)
będzie odwzorowaniem \(T(x)=Ax\).
Jego odwzorowanie dualne ma
w bazach dualnych macierz
\(A^{\mathsf T}\) na mocy
Twierdzenia~\ref{thm:lin7-macierz-dualnego}.
Wynik wynika z
Twierdzenia~\ref{thm:lin7-jadro-obraz-dualnego}.
\end{proof}

\begin{wniosek}[Wzajemność iniektywności i surjektywności]
\label{cor:lin7-injekcja-surjekcja}
Niech \(V,W\) będą skończenie
wymiarowymi przestrzeniami,
a \(T\colon V\to W\) liniowym
odwzorowaniem. Wtedy
\[
 \begin{aligned}
 T\text{ jest surjektywne}
 &\quad\Longleftrightarrow\quad
 T^*\text{ jest iniektywne},\\
 T\text{ jest iniektywne}
 &\quad\Longleftrightarrow\quad
 T^*\text{ jest surjektywne}.
 \end{aligned}
\]
\end{wniosek}
\begin{proof}
Pierwszą równoważność otrzymujemy,
obliczając
\(\ker T^*=(\operatorname{im}T)^\circ\).
Anihilator podprzestrzeni
\(M\subset W\) jest zerowy dokładnie
wtedy, gdy \(M=W\), ponieważ
\(\dim M^\circ=\dim W-\dim M\).
Zatem \(\ker T^*=\{0\}\)
dokładnie wtedy, gdy
\(\operatorname{im}T=W\).

Drugą równoważność otrzymujemy z
\(\operatorname{im}T^*=(\ker T)^\circ\).
Anihilator \(\ker T\) jest całym \(V^*\)
wtedy i tylko wtedy, gdy \(\ker T=\{0\}\),
co można sprawdzić porównując wymiary.
\end{proof}

\section{Przestrzeń bidualna}
\label{sec:lin7-bidualnosc}

\begin{definicja}[Przestrzeń bidualna i odwzorowanie kanoniczne]
\label{def:lin7-bidual}
Przestrzenią \emph{bidualną} do \(V\)
nazywamy \(V^{**}:=(V^*)^*\).
Każdemu \(v\in V\) odpowiada
funkcjonał \(J_V(v)\in V^{**}\)
zdefiniowany przez
\[
 J_V(v)(\varphi):=\varphi(v)
 \qquad(\varphi\in V^*).
\]
Otrzymujemy odwzorowanie
\(J_V\colon V\to V^{**}\)
nazywane \emph{odwzorowaniem
kanonicznym do przestrzeni bidualnej}.
\end{definicja}

\begin{twierdzenie}[Bidualność w wymiarze skończonym]
\label{thm:lin7-bidualnosc}
Niech \(V\) będzie skończenie wymiarowa.
Wtedy \(J_V\colon V\to V^{**}\)
jest kanonicznym izomorfizmem
przestrzeni wektorowych.
Jeżeli \(B=(v_1,\ldots,v_n)\)
jest uporządkowaną bazą \(V\),
to \((J_V(v_1),\ldots,J_V(v_n))\)
jest bazą dualną do \(B^*\)
w przestrzeni \(V^{**}\).
\end{twierdzenie}
\begin{proof}
Dla każdego \(v\in V\)
funkcja \(\varphi\mapsto\varphi(v)\)
jest liniowa z \(V^*\) do \(K\),
zatem \(J_V(v)\in V^{**}\).
Dla \(a,b\in K\), \(v,w\in V\)
i dowolnego \(\varphi\in V^*\)
mamy
\[
 J_V(av+bw)(\varphi)
 =\varphi(av+bw)
 =aJ_V(v)(\varphi)
  +bJ_V(w)(\varphi).
\]
Stąd \(J_V\) jest liniowe.

Jeśli \(J_V(v)=0\), to
\(\varphi(v)=0_K\) dla każdego
\(\varphi\in V^*\).
Ze
Stwierdzenia~\ref{prop:lin7-ewaluacja}
wynika \(v=0\), więc \(J_V\)
jest iniektywne.
Z Twierdzenia~\ref{thm:lin7-baza-dualna}
otrzymujemy
\(\dim V^{**}=\dim V^*=\dim V\).
Liniowe odwzorowanie iniektywne
między skończenie wymiarowymi
przestrzeniami równego wymiaru
jest surjektywne, a więc \(J_V\)
jest izomorfizmem.

Ponadto
\(J_V(v_i)(v^j)=v^j(v_i)=\delta_{ij}\),
czyli obrazy wektorów bazy \(B\)
są elementami bazy dualnej
do \(B^*\). Dotyczy to także
przypadku \(n=0\).
\end{proof}

\begin{twierdzenie}[Naturalność identyfikacji z bidualem]
\label{thm:lin7-naturalnosc-bidualu}
Niech \(V,W\) będą skończenie
wymiarowymi przestrzeniami
i \(T\colon V\to W\) odwzorowaniem
liniowym. Wtedy
\[
 T^{**}\circ J_V=J_W\circ T,
\]
gdzie \(T^{**}:=(T^*)^*
\colon V^{**}\to W^{**}\).
W szczególności utożsamienie
\(V\cong V^{**}\) nie zależy
od wyboru bazy.
\end{twierdzenie}
\begin{proof}
Niech \(v\in V\) i \(\psi\in W^*\).
Z definicji odwzorowań dualnych
otrzymujemy
\[
 \begin{aligned}
 (T^{**}(J_Vv))(\psi)
 &=J_Vv(T^*\psi)\\
 &=(T^*\psi)(v)
 =\psi(Tv)
 =J_W(Tv)(\psi).
 \end{aligned}
\]
Oba elementy \(W^{**}\) mają
te same wartości na każdym
\(\psi\in W^*\), więc są równe.
Formuła definiująca \(J_V\)
nie wykorzystuje bazy;
ostatni wniosek jest więc
własnością samej konstrukcji.
\end{proof}

\begin{uwaga}[Wymiar nieskończony]
\label{rem:lin7-nieskonczony-bidual}
Dla dowolnej przestrzeni wektorowej
odwzorowanie \(J_V\) ma sens i jest
liniowe. Przy istnieniu baz Hamela
pozostaje iniektywne: niezerowy wektor
można włączyć do bazy i przypisać mu
wartość \(1_K\) funkcjonału liniowego.
W przestrzeniach nieskończenie
wymiarowych \(J_V\) nie musi być
surjektywne. Dlatego twierdzenie
o izomorfizmie sformułowaliśmy
wyłącznie w wymiarze skończonym.
\end{uwaga}

\section{Ślad endomorfizmu}
\label{sec:lin7-slad}

\begin{definicja}[Ślad macierzy]
\label{def:lin7-slad}
Niech \(n\in\mathbb N_0\) i
\(A=(a_{ij})\in\operatorname{Mat}_{n\times n}(K)\).
\emph{Śladem} \(A\) nazywamy
\[
 \operatorname{tr}A:=\sum_{i=1}^n a_{ii}.
\]
Dla \(n=0\) jest to pusta suma,
czyli \(\operatorname{tr}I_0=0_K\).
\end{definicja}

\begin{twierdzenie}[Liniowość i cykliczność śladu]
\label{thm:lin7-slad}
Niech \(n\in\mathbb N_0\).
Dla dowolnych macierzy
\(A,B\in\operatorname{Mat}_{n\times n}(K)\)
i \(a,b\in K\) mamy
\[
 \operatorname{tr}(aA+bB)
 =a\operatorname{tr}A+b\operatorname{tr}B,
 \qquad
 \operatorname{tr}(AB)
 =\operatorname{tr}(BA).
\]
Jeśli \(P\in\operatorname{Mat}_{n\times n}(K)\)
jest odwracalna, to
\[
 \operatorname{tr}(P^{-1}AP)
 =\operatorname{tr}A.
\]
Dlatego dla dowolnego endomorfizmu
\(T\colon V\to V\) skończenie
wymiarowej przestrzeni \(V\)
można poprawnie zdefiniować
\(\operatorname{tr}T\)
jako ślad macierzy \(T\)
w dowolnej uporządkowanej bazie.
\end{twierdzenie}
\begin{proof}
Liniowość wynika wprost ze wzoru
na ślad i punktowych działań
na macierzach. Dla cykliczności
rozpisujemy wyrazy diagonalne:
\[
 \begin{aligned}
 \operatorname{tr}(AB)
 &=\sum_{i=1}^n\sum_{j=1}^n a_{ij}b_{ji}\\
 &=\sum_{j=1}^n\sum_{i=1}^n b_{ji}a_{ij}
 =\operatorname{tr}(BA).
 \end{aligned}
\]
W drugim przejściu użyliśmy
przemienności mnożenia w ciele \(K\)
oraz możliwości zmiany kolejności
skończonych sum.

Dla odwracalnej macierzy \(P\)
cykliczność zastosowana do
\(P^{-1}A\) i \(P\) daje
\[
 \operatorname{tr}(P^{-1}AP)
 =\operatorname{tr}(PP^{-1}A)
 =\operatorname{tr}A.
\]
Macierze jednego endomorfizmu
w różnych bazach są podobne
na mocy
Wniosku~\ref{cor:lin4-podobienstwo},
więc ich ślady są jednakowe.
W wymiarze \(0\) wszystkie sumy
w dowodzie są puste; w szczególności
ślad jest równy \(0_K\).
\end{proof}

\begin{stwierdzenie}[Ślad odwzorowania rzędu co najwyżej jeden]
\label{prop:lin7-slad-rank-jeden}
Niech \(V\) będzie skończenie wymiarowa,
\(v\in V\) oraz \(\varphi\in V^*\).
Zdefiniujmy endomorfizm
\(R_{v,\varphi}\colon V\to V\) wzorem
\[
 R_{v,\varphi}(x):=\varphi(x)v.
\]
Wtedy \(\operatorname{rk}R_{v,\varphi}\leq1\)
i
\[
 \operatorname{tr}R_{v,\varphi}=\varphi(v).
\]
\end{stwierdzenie}
\begin{proof}
Obraz \(R_{v,\varphi}\)
zawiera się w \(\operatorname{span}\{v\}\),
więc jego wymiar jest co najwyżej \(1\).
Niech \(B=(v_1,\ldots,v_n)\)
będzie dowolną uporządkowaną bazą,
a \([v]_B=(c_1,\ldots,c_n)^{\mathsf T}\).
Ponieważ
\(R_{v,\varphi}(v_j)=\varphi(v_j)v\),
element \((i,j)\) macierzy
tego endomorfizmu jest równy
\(c_i\varphi(v_j)\).
Zatem jej ślad wynosi
\[
 \sum_{i=1}^n c_i\varphi(v_i)
 =\varphi\left(\sum_{i=1}^n c_iv_i\right)
 =\varphi(v).
\]
Dla \(V=\{0\}\) mamy \(v=0\),
funkcjonał zerowy i ślad \(0_K\),
więc wzór nadal zachodzi.
\end{proof}

\begin{przyklad}[Ślad a wyznacznik]
\label{prz:lin7-slad-wyznacznik}
Dla \(A=\begin{pmatrix}1&2\\0&3\end{pmatrix}\)
nad \(\mathbb R\) mamy
\(\operatorname{tr}A=1+3=4\),
a \(\det A=1\cdot3=3\).
Są to różne niezmienniki
podobieństwa. Dla
\(B=\begin{pmatrix}3&0\\0&1\end{pmatrix}\)
otrzymujemy ten sam ślad i
ten sam wyznacznik, lecz sama zgodność
tych dwóch liczb nie dowodzi podobieństwa
dowolnych macierzy.
\end{przyklad}

\section{Podsumowanie}
\label{sec:lin7-podsumowanie}

Przestrzeń dualna \(V^*\)
składa się ze wszystkich funkcjonałów
liniowych na \(V\). W wymiarze
skończonym każdej bazie \(V\)
odpowiada baza dualna, a współrzędne
funkcjonału są jego wartościami
na wektorach wyjściowej bazy.
Anihilator \(U^\circ\)
zawiera dokładnie funkcjonały
zerujące się na \(U\), przy czym
\(\dim U^\circ=\dim V-\dim U\).
Odwzorowaniu liniowemu \(T\colon V\to W\)
odpowiada odwzorowanie dualne
\(T^*\colon W^*\to V^*\)
o macierzy transponowanej.
Kanoniczny izomorfizm
\(V\to V^{**}\) istnieje
w wymiarze skończonym bez wyboru bazy;
nie należy go mylić z zależnym
od bazy izomorfizmem \(V\cong V^*\).
Ślad endomorfizmu, podobnie jak
wyznacznik, jest niezmiennikiem
zmiany bazy, lecz opisuje inną
własność odwzorowania.
